% ===================================================================
% Back cover
% ===================================================================
\clearpage
\thispagestyle{empty}
\null\vfill

\begin{center}
{\color{ruleblue}\rule{0.5\textwidth}{1pt}}\par
\vspace{0.4cm}
{\normalfont\Large\bfseries Lean for Working Algebraists\par}
\vspace{0.4cm}
{\color{ruleblue}\rule{0.5\textwidth}{1pt}}
\end{center}

\vspace{0.8cm}
\begin{quote}
\noindent
A guide to Lean~4 for readers who already think in groups, rings,
functors, and diagrams, but have never written a line of Lean (or any
proof assistant). No programming background is assumed --- only the
mathematical maturity of someone comfortable with abstract algebra and
the basic language of category theory. Starting from \texttt{\#eval} and
\texttt{def}, this book builds groups, rings, modules, and quiver path
algebras entirely from scratch, then hands off to Mathlib once the
underlying ideas are firmly in hand.

\medskip
\noindent\textbf{Inside:}
\begin{itemize}
\item Dependent types, $\Pi$/$\Sigma$-types, and the calculus of
  constructions, built from concrete \texttt{Fin}/\texttt{Vec} examples
  rather than abstract nonsense.
\item Groups, rings, modules, and quiver path algebras, formalized by
  hand, field by field, with every proof obligation spelled out.
\item Two checkpoint projects and five scaffolded next-step projects,
  each with learning objectives, milestones, and a self-verification
  step.
\item A consolidated bibliography, a tactic and Mathlib-name reference,
  and worked solutions to every exercise.
\end{itemize}

\medskip
\noindent Every Lean code block in this book was checked against a real
\texttt{v4.31.0} toolchain, not merely described.
\end{quote}

\vspace{0.8cm}
\begin{center}
{\normalfont\itshape Written by Abderrahim Adrabi, in collaboration with
Claude Code.\par}
\end{center}

\vfill\null
